\scanpage{51}
\mspage{5}

\uline{Proposition 1} Soit $X$ une catégorie telle que

\deleted{\unsure{a)}} a) $X$ pseudocofiltrante \quad b) $X$ admet un plus petit objet $e$ \inserted{(conséquence de \unsure{a)} si $\unsure{\operatorname{Ord}|\operatorname{Ob}} X$ fini)} \deleted{objet minimal \unsure{$e$}} c) Pour tt $x\in\operatorname{Ob}X$, l'ens des sous-$\underbrace{\operatorname{Hom}(e,e)}_{M}$-ens. à droite de $\operatorname{Hom}(e,x)$ satisfait la condition minimale (ce qui est le cas si les $\operatorname{Hom}(e,x)$ sont finis). Alors

$1^\circ)$ $\forall x\in\operatorname{Ob}X$, parmi les sous-$M$-ens. à gauche de $\operatorname{Hom}(e,x)\overset{\mathrm{def}}=F(x)$ qui sont de la forme $fM$ (avec $f\in F(x)$), il y en a un plus petit, \inserted{soit $H_x=\bigcap_{f\in F(x)}fM$.} Il est formé de l'\unsure{ens|un} des $\lambda_{\deleted{\unsure{x}}}\in\operatorname{Hom}(e,x)$ \deleted{tels que} \inserted{qui sont} tels que \deleted{$\forall f\in\operatorname{Hom}($} $\forall f\in F(x)$, $\exists u\in M$ avec $\lambda=fu$.

$2^\circ)$ $H_x$ est donc un $M$-ens. à droite. De plus, pour tout $\deleted{\unsure{x}}f:x\longrightarrow y$ dans $X$, on a $H_y=f.H_x$. On a donc $\operatorname{card} H_y\leq\operatorname{card} H_x$ si $x\leq y$ ($\overset{\mathrm{def}}\Longleftrightarrow$ si $\operatorname{Hom}(x,y)\ne\varnothing$).

$3^\circ)$ Supposons les $F(x)$ finis (cf condition c) ci-dessus) et que $e$ soit aussi objet maximal, i.e. que $\operatorname{Ord} X$ soit réduit
% archive sheet number: 50
